\section{Transfer matrices and growth shortfalls}
\label{sec:setup}

Throughout the proof the parameter $k$ is positive and sufficiently large.
Set
\begin{equation}\label{setup:parameters}
 K=2\pi k,\qquad M=K+4,\qquad
 \phi(x)=2x+k\sin(2\pi x),\qquad
 v(x)=\phi'(x)=2+K\cos(2\pi x).
\end{equation}
Growth logarithms are to base $M$. Natural logarithms are retained in the
definitions of entropy and Lyapunov exponents. Constants denoted by $C$
may increase from one occurrence to the next. Unless a dependence is
specified, constants in large-$M$ estimates do not depend on the orbit,
the integer times, or $M$.

Write $(x_i,y_i)=f_k^i(z)$ for all $i\in\Z$, and put $q_i=x_i\in\T$.
The identity $x_i-y_i=x_{i-1}$ shows that the fixed torus automorphism
$(x,y)\mapsto(x,x-y)$ replaces the state by $(q_i,q_{i-1})$.
It preserves normalized area, and the recurrence becomes
\begin{equation}\label{setup:recurrence}
 q_{i+1}=\phi(q_i)-q_{i-1}\pmod1.
\end{equation}
When varying an orbit locally we choose initial real lifts and iterate
this same recurrence over $\R$, without reducing modulo one. Its linear
recurrence is
\begin{equation}\label{setup:transfer}
 \binom{u_{i+1}}{u_i}
 =A_i\binom{u_i}{u_{i-1}},\qquad
 A_i=\begin{pmatrix}v(q_i)&-1\\1&0\end{pmatrix}.
\end{equation}
For $a,b\in\Z$, let $T_{a,b}$ carry the linearized pair at time $a$ to
the pair at time $b$. Thus $T_{a,c}=T_{b,c}T_{a,b}$ and
$T_{b,a}=T_{a,b}^{-1}$.

All matrix and vector norms are Euclidean unless otherwise stated.
Every transfer matrix has determinant one. A real two-dimensional
determinant-one matrix has singular values $D,D^{-1}$ with $D\ge1$;
in particular
\begin{equation}\label{setup:norm-bounds}
 \norm{T_{a,b}}=\norm{T_{b,a}},\qquad
 1\le\norm{T_{a,b}}\le M^{|b-a|}.
\end{equation}
Indeed both $A_i$ and $A_i^{-1}$ have Euclidean norm at most $M$.
The same bound holds in the maximum norm: their largest absolute row sum
is at most $K+3<M$. Also
\begin{equation}\label{setup:derivatives}
 |\phi'|\le M,\qquad |\phi^{(p)}|\le C_pM\quad(p\ge2)
\end{equation}
for every fixed derivative order $p$. These estimates apply equally on
all real lifts. The fixed change of coordinates does not alter Lyapunov
exponents, since conjugating a derivative product by a fixed invertible
matrix changes its logarithmic norm by a bounded amount.

The probability $\PP$ and expectation $\E$ always refer to normalized
area of the initial state. At every fixed time $i$, the pair
$(q_i,q_{i-1})$ has uniform product distribution on $\T^2$.
The process is stationary under all integer time shifts. This does not
assert independence of coordinates observed at different times.

We record growth and its expected shortfall by
\begin{equation}\label{scale:distance-deficit}
 g(i,j)=\log_M\norm{T_{i,j}},\qquad
 e_n=\E\left[1-\frac{g(0,n)}n\right]\quad(n\ge1).
\end{equation}
The transfer identities make $g$ a stationary semimetric on orbit time:
it is symmetric, vanishes on the diagonal, and satisfies the triangle
inequality. Distinct times may have zero distance. The one-step bound gives
\begin{equation}\label{scale:distance-bound}
 0\le g(i,j)\le|i-j|,
\end{equation}
so $e_n\in[0,1]$ measures the average loss from the upper growth rate
over an interval of length $n$. These definitions do not assume that
Lyapunov exponents vanish; that assumption enters only in the
contradiction argument of \cref{sec:scales}.

For nonzero vectors $a,b\in\R^2$, the symbol $a\wedge b$ denotes their
determinant. For unit vectors its absolute value is the absolute sine
of their angle and depends only on their unoriented lines. The absolute
sines satisfy
\begin{equation}\label{setup:sine-triangle}
 |a\wedge c|\le |a\wedge b|+|b\wedge c|
\end{equation}
for unit vectors: express the two outer vectors in the orthonormal basis
given by $b$ and a perpendicular vector, and expand their determinant.

The following elementary observation will repeatedly convert a small
total growth deficit into good growth on many subintervals.

\begin{lemma}[Shortfall count]\label{setup:shortfall-count}
Let $z_0,\ldots,z_N$ be real numbers with $z_j-z_{j-1}\le1$, and set
\[
 \mathcal D=N-(z_N-z_0)
 =\sum_{j=1}^N\bigl(1-z_j+z_{j-1}\bigr).
\]
For $c<1$, the number of indices $j$ for which some $0\le l<j$ satisfies
$z_j-z_l<c(j-l)$ is at most $C\mathcal D/(1-c)$, with an absolute
constant $C$. The same conclusion holds for the number of indices $j$
with some $j<l\le N$ satisfying $z_l-z_j<c(l-j)$.
\end{lemma}

\begin{proof}
For each bad index choose one witnessing interval. Its length is a
positive integer, and the sum of the nonnegative step shortfalls over
it exceeds $(1-c)$ times its length. Select disjoint intervals greedily
in decreasing order of length, discarding each interval that intersects
one already selected. Regard intervals as half-open when deciding
disjointness, so that the sums of shortfalls over selected intervals
have disjoint sets of steps.

Every discarded interval meets a selected interval of at least its
length and is contained in the concentric threefold enlargement of that
selected interval. The bad endpoints are therefore covered by these
enlargements. An interval of integer length $L\ge1$ has at most $3L+1$
integer points in its threefold enlargement. If the selected lengths
are $L_1,\ldots,L_s$, the number of bad indices is at most
$4\sum_r L_r$. Disjointness and nonnegativity give
$(1-c)\sum_rL_r\le\mathcal D$. This proves the first statement, and
the same selection applied to the forward witnessing intervals proves
the second.
\end{proof}

We will use the count both for logarithms of vector norms, whose
one-step increments are at most one by \cref{setup:norm-bounds}, and
for the logarithms of successive maximum norms of solutions of
\cref{setup:transfer}. Neither use requires a lower bound on each
individual increment.
